% !TEX root = ../monograph.tex
\section{Calibrating the microscopic equations on the same history}
\label{ch11:equation-calibration}

Fix a degree $d$. The section $\mathcal V_d$ and automorphism $R_d$
are those of Theorem~\ref{ch02:equation-adapted}; neither is chosen
again when the cutoff or period changes. The physical section remains
$P_d^{\rm loc}$. We compare the two on the history already used to
define $\mathsf Z_{N,p,d}$.

We first specify an additional finite stage requirement.
At stage $j$, let $a_{dj}$ bound the derivatives of $\mathcal V_d$
requested by that stage, for every $d\le j$. These are finite constants
by the powered-curvature-factor estimate in Chapter~\ref{ch02:observables}.
Let $q_d$ bound the fixed union stencil and put
\[
 Q_{dj}=2+e^{81\sigma_j}(a_{dj}+C_j+2q_d).
\]
Here $\sigma_j$ is the support exponent and $C_j$ the column bound.
These quantities are fixed before the block length.
Let $\mathcal O_d\ge1$ bound both the finite symmetry actions on
coefficient $\ell^1$ norm and their averaging coefficients, and put
$\delta_d=(2\mathcal O_d^2)^{-1}$. Include the finite maximum
\begin{equation}
 L^{1/8}\ge
 \max_{d\le j}2(1+C_j)Q_{dj}\delta_d^{-1}
 \label{ch11:calibration-floor}
\end{equation}
in the stage-$j$ block floor. Include
$\widehat\beta_N\ge a_{d,J(N)}$ in the adjacent stage conditions on the starting scale. These requests are compatible with the selection order:
the section and its derivative constants are fixed first, then the
length, then the inverse coupling threshold. Raising a later threshold does
not change the section. Thus they can be imposed for all stages
before selecting the common history, with no dependence on $p$.

Define the normalized microscopic section at the origin by
\[
 O^{\mathcal V}_{N,p}(v)
   =a_N^{-d}\mathcal V_d[S_{1/t_N}v](0).
\]
The factor scaling $S_{1/t_N}$ commutes with $R_d$ only when
$R_d$ preserves factor count, which is not asserted here.
Accordingly its exact equation labels at this normalization are
\begin{equation}
 I_{N,d}^{\mathcal V}
    =S_{t_N}R_d\mathcal I_d.
 \label{ch11:scaled-equation-labels}
\end{equation}
Since $\mathcal I_d$ is homogeneous in factor count, one can equally
write $I_{N,d}^{\mathcal V}=
(S_{t_N}R_dS_{1/t_N})\mathcal I_d$.
This conjugation must be retained when using a cutoff-independent
unscaled equation section.

\begin{proposition}[One comparison map for all equation labels]
\label{ch11:equation-comparison}
Under \eqref{ch11:calibration-floor} and the terminal attachment
bounds of Section~\ref{ch11:trial-descent}, there is, for every
sufficiently large $N$ and every $p$, one real equivariant map
$P_{N,p,d}:V_d\to V_d$ such that
\[
 \|P_{N,p,d}-I\|_{1\to1}\le\tfrac12,\qquad
 \|P_{N,p,d}^{-1}\|_{1\to1}\le2.
\]
For each fixed stopping pair $(m,k)$,
\begin{equation}
 \mathcal R_m\!\left\{
 O^{\mathcal V}_{N,p}(v)
   -O_{N,p}(d,\mathsf Z_{N,p,d}P_{N,p,d}v;0)\right\}
 \le C_{d,m,k}\|v\|_1 E_N e^{-\tau_N/8},
 \label{ch11:equation-root-error}
\end{equation}
where $E_N$ is subpower in $a_N^{-1}$, uniformly in $p$.
The same map works for every spectator; it is constructed before
any spectator is specified.
\end{proposition}
\begin{proof}
For each basis vector, the exact terminal decomposition used in
Section~\ref{ch11:trial-descent} applies to the fixed section
$\mathcal V_d$. Its part with zero polynomial coefficient is localized on the union
stencil; the complementary part is attached to the original
large-field selection, with no change of density.
It gives coefficient $e_j$, ordinary norm bound
$Q_N=Q_{d,J(N)}$, and complete source norm bound
\[
 E_N=4(1+C_N)Q_N+
 e^{81\sigma_N+162}Z_N^{\rm val}\zeta_{d,N}u_N^{-1}.
\]
Here
$Z_N^{\rm val}=(A_{\mathcal V}+2A_{\rm loc})
 \sum_jt_N^{-s_j}$, with $A_{\mathcal V},A_{\rm loc}$
the global value norms of the two fixed sections.
The ordinary terminal norm bound uses a sharper attachment estimate
than the general trial bound. If $w_N^{\rm for}$ is the terminal
large-field bound and $D$ the union stencil, its contribution
to the remainder norm normalized by $\sqrt{w_N^{\rm for}}$ is
\[
 2|D|\sqrt{w_N^{\rm for}}\,
       e^{\sigma_N c_{B_N}(D)}Z_N^{\rm val}.
\]
Indeed attach the complementary local function to the original
nonempty prescribed large-field set and sum its possible distinguished
sites in $D$; the remaining large-field suppression bounds this
normalized mass. The same exact attachment retains reserved mass
$e^{(\sigma_N+2)c_{B_N}(D)}Z_N^{\rm val}$.
Impose the additional conditions on the starting scale
\[
 |D|w_N^{\rm for}\le1,\qquad c_{B_N}(D)\le81,\qquad
 \sqrt{w_N^{\rm for}}Z_N^{\rm val}\le1.
\]
Then $|D|\le q_d$ bounds the ordinary contribution by
$2q_de^{81\sigma_N}$, explaining exactly the definition of $Q_N$.
These are post-length threshold conditions. At fixed adjacent
stages the value norm is a fixed polynomial in the inverse
small-field threshold, whereas the large-field bound decays
faster than every such power. Their finite simultaneous threshold
therefore exists. The unweighted reserved mass is not discarded.
This last term is required: the ordinary derivative bound alone
does not bound the complementary large-field contribution.

Backward comparison gives a reference coefficient $v_j$ and the
chessboard-seminorm error in \eqref{ch11:trial-error}. Put
$u_j=\mathsf Z_{N,p,d}^{-1}v_j$.
There is a slightly stronger endpoint estimate than the displayed
bulk descent bound. At the last step its cone inequality gives
\[
 |e_j-u_j|\le
 4K_{N-1}L_{N-1}^{-3/8}c_N
 \le L_{N-1}^{-1/8}c_N
 \le 2(1+C_N)Q_N L_{N-1}^{-1/8}.
\]
The middle inequality follows from the previously imposed block
floor. Requirement \eqref{ch11:calibration-floor} bounds this by
$\delta_d$.

Let $Ue_j=u_j$. Average $U$ by conjugation with the finite geometric
symmetry group, and then with the real structure, obtaining $P$.
For the conjugation term use $\overline{U\overline v}$, so the result
is still complex linear. The norm of $P-I$ is at most
$\mathcal O_d^2\delta_d=1/2$, and the Neumann series gives the stated
inverse bound. Apply the same averaging to the field errors.
Linearity and the triangle inequality for the chessboard seminorm give \eqref{ch11:equation-root-error}, enlarging its fixed
constant to cover the finite orbit. This preserves one common
history, all source monomials, and one common map.

Finally $C_N$, $e^{81\sigma_N}$, $t_N^{-1}$,
$\zeta_{d,N}$, and $u_N^{-1}$ have the subpower bounds already
required for trial descent. The extra condition on the starting scale bounds
$a_{d,J(N)}$ by $\widehat\beta_N$, which is subpower.
Finite sums and products preserve this property, proving the claim
for $E_N$.
\end{proof}

The varying equation subspace in reference coordinates is therefore
\begin{equation}
 \mathsf E_{N,p,d}
   =\mathsf Z_{N,p,d}P_{N,p,d}
               S_{t_N}R_d\mathcal I_d.
 \label{ch11:actual-equation-subspace}
\end{equation}
For $v_N\in I_{N,d}^{\mathcal V}$ of subpower norm, the first term
in \eqref{ch11:equation-root-error} has exactly zero correlation
against every bounded spectator outside its microscopic stencil,
by \eqref{ch02:microscopic-null}. For fixed positive physical
separation, this exclusion holds at every sufficiently fine mesh,
uniformly in period. Holder's chessboard bound then shows that the
physical insertion with coefficient
$\mathsf Z_{N,p,d}P_{N,p,d}v_N$ has vanishing separated correlations
against any finite family of spectators with bounded chessboard seminorms. This proves the equation inclusion needed for the
argument using a finite separating family; equality with the full null space still
uses the free nondegeneracy and detection estimates.

The scaling in \eqref{ch11:actual-equation-subspace} also has a
useful check. Since $R_d-I$ never lowers factor count and $t_N\to0$,
the conjugate $S_{t_N}R_dS_{1/t_N}$ remains bounded.
Its strictly factor-raising terms tend to zero. Terms at equal
factor count can still raise scaling dimension and need not tend
to zero under this scaling alone. Thus the raw triangular
construction and the near-identity comparison $P$ do not, by
themselves, prove that this equation subspace approaches the
unmodified Maxwell ideal. No such identification is required here: use bounded projections
with kernel exactly \eqref{ch11:actual-equation-subspace} and establish
the bound from the separating family \eqref{ch11:detector-bound} on their complements.
The terminal map $\mathsf Z$ already includes the scaling-dimension
evolution; its physical convention has not been changed.
